\documentclass[10pt,a4paper]{amsart}
\usepackage[margin=19mm]{geometry}
\usepackage{amsmath,amssymb,mathtools}
\usepackage[scr=rsfs,cal=euler]{mathalfa}
\usepackage{xfrac}
\usepackage[unicode,hidelinks,bookmarks=false]{hyperref}
\newcommand{\ie}{i.e.\ }
\newcommand{\cf}{cf.\ }
\newcommand{\resp}{respectively\ }
\newcommand{\Subsection}{\subsection}
\providecommand{\nobreakdash}{\nobreak\hbox{-}}
\setlength{\emergencystretch}{2em}
\multlinegap0pt
\usepackage{amssymb}
\newtheorem{lemma}{Lemma}
\newtheorem{theorem}{Theorem}
\newtheorem{proposition}{Proposition}
\title{Rigidity of Euclidean Minimal Hypersurfaces under Nonuniform Diagonal Dilations}
\author{Jongha Lee, Suhwan Lee, Jae Won Lee}
\date{Source: arXiv:2609.00668v1 (CC BY 4.0)}
\begin{document}
\maketitle

\noindent\emph{Source and licence.} Rigidity of Euclidean Minimal Hypersurfaces under Nonuniform Diagonal Dilations. Version arXiv:2609.00668v1 (CC BY 4.0), distributed by arXiv under the Creative Commons Attribution 4.0 International licence, http://creativecommons.org/licenses/by/4.0/, as recorded in the arXiv licence metadata for this exact version and shown on the abstract page https://arxiv.org/abs/2609.00668v1. Full licence notice: \texttt{licenses/arxiv-2609.00668-CC-BY-4.0.txt}.\par\smallskip

\noindent\textit{Editorial scope.} Counted unit: the article's proposition prop:repeatedweights (source0.tex lines 350--352). The article's complete original proof of it is delivered with it (source0.tex lines 353--393, one \texttt{proof} environment of 41 lines). No mathematics is written, repaired or re-derived: the statement and the proof are the article's own lines, in the article's own notation, and no retained line is substituted or re-flowed.  Retained uncounted as the source-local context the proof needs: lines 146--148; lines 173--186; lines 175--177; lines 264--271; lines 313--327.  Nothing was omitted from any retained span, so the package carries no omission record.  No retained line contains a citation, so the wrapper carries no bibliography and no bibliography entry.  No retained line contains a figure environment, an image-inclusion command or drawing code, so no image is delivered and the package declares no asset dependency.  Editorial formatting only, in addition: an AMS \texttt{amsart} preamble that restates the article's own declarations and macros used by the retained lines, namely the environments \texttt{lemma}, \texttt{theorem}, \texttt{proposition} on separate counters, together with the standard packages those lines need.  The retained environments print in this document as Lemma 1, Theorem 1, Proposition 1 in order of appearance, whereas the article prints the same items with the numbering of its own full document.  The complete native source of the article as distributed in the arXiv e-print for this exact version is bundled unchanged as \texttt{sources/209048/source0.tex}.
\begin{equation}\label{eq:dilationidentity}
\mathcal M[F_t](D_tx)=\sum_{1\le i<j\le n}t^{-2(g_i+g_j)}C_{ij}[F](x).
\end{equation}

\begin{lemma}[pair-vector identities]\label{lem:pairvectors}
Along $\Sigma$, every $V_{ij}$ is tangent and
\begin{equation}\label{eq:CijII}
C_{ij}[F]=\operatorname{Hess}F(V_{ij},V_{ij})=|\nabla F|\,II(V_{ij},V_{ij}),
\end{equation}
up to orientation. Moreover,
\begin{equation}\label{eq:traceidentity}
\sum_{i<j}V_{ij}V_{ij}^{\!\top}=|\nabla F|^2I-\nabla F\nabla F^{\!\top},
\end{equation}
and consequently
\begin{equation}\label{eq:minop-pairs}
\mathcal M[F]=\sum_{i<j}\operatorname{Hess}F(V_{ij},V_{ij}).
\end{equation}
\end{lemma}

\begin{equation}
C_{ij}[F]=\operatorname{Hess}F(V_{ij},V_{ij})=|\nabla F|\,II(V_{ij},V_{ij}),
\end{equation}

\begin{theorem}[finite-dilation rigidity and affine characterization]\label{thm:main}
Let $n\ge3$, let $\Sigma\subset\mathbb R^n$ be a connected embedded $C^2$ hypersurface, and let $D_t$ have positive pair-sum-nonresonant weights. Put $N=\binom n2$. The following are equivalent:
\begin{enumerate}
\item $D_t\Sigma$ is minimal for every $t>0$;
\item $D_{t_1}\Sigma,\ldots,D_{t_N}\Sigma$ are minimal for some $N$ distinct positive numbers $t_1,\ldots,t_N$;
\item $\Sigma$ is a connected open subset of an affine hyperplane.
\end{enumerate}
\end{theorem}

\begin{proposition}[nonflat minimal geometry in resonant regimes]\label{prop:resonance}
Pair-sum resonance can preserve nonflat minimal geometry under genuinely nonuniform diagonal dilations.
\begin{enumerate}
\item In $\mathbb R^3$, let $g_1=g_2=1$, $g_3=\gamma>0$, $\gamma\ne1$, and on the quadrant $x_1,x_2>0$ set
\begin{equation}\label{eq:helicoid}
\Sigma_c=\{x_3=c\arctan(x_2/x_1)\},\qquad c\ne0.
\end{equation}
Then every $D_t\Sigma_c$ is again a helicoid patch and hence minimal, while $\Sigma_c$ is nonflat.
\item In $\mathbb R^4$, let $g=(1,2,3,4)$ and
\begin{equation}\label{eq:cone}
\Sigma=\{x_1x_4+x_2x_3=0\}.
\end{equation}
Its regular part is $\Sigma\setminus\{0\}$; it is connected, nonflat, and minimal, and it satisfies $D_t(\Sigma\setminus\{0\})=\Sigma\setminus\{0\}$ for every $t>0$.
\end{enumerate}
\end{proposition}

\begin{proposition}[repeated weights destroy rigidity in every dimension]\label{prop:repeatedweights}
Let $n\ge3$. If two dilation weights coincide, then there exists a connected real-analytic embedded nonflat minimal hypersurface patch whose image under every $D_t$ is again minimal. Thus, apart from pair-sum nonresonance, the example satisfies the geometric hypotheses appearing in Theorem~\ref{thm:main}; repeated weights are a structural resonant obstruction to its rigidity conclusion.
\end{proposition}

\begin{proof}
After relabeling, suppose $g_1=g_2$. For $n=3$ use the helicoid patch in Proposition~\ref{prop:resonance}(1). For $n\ge4$, choose constants $a_3,\ldots,a_{n-1}$ and $c\ne0$ and consider, on the quadrant $x_1,x_2>0$, the graph
\begin{equation}\label{eq:shearedhelicoid}
x_n=a_3x_3+\cdots+a_{n-1}x_{n-1}+c\arctan(x_2/x_1).
\end{equation}
Write $h=L+c\theta$, where $L=a_3x_3+\cdots+a_{n-1}x_{n-1}$ and $\theta=\arctan(x_2/x_1)$. Since $\theta$ is harmonic and $\operatorname{Hess}\theta(\nabla\theta,\nabla\theta)=0$, the graph minimal-surface operator satisfies
\[
(1+|\nabla h|^2)\Delta h-\operatorname{Hess}h(\nabla h,\nabla h)=0.
\]
Thus the graph is minimal. Under $D_t$, write $y=D_tx$. Solving for the last coordinate gives
\[
y_n=\sum_{k=3}^{n-1}a_k t^{g_n-g_k}y_k+c\,t^{g_n}\arctan\!\left(\frac{y_2}{y_1}\right),
\]
because
\[
\frac{x_2}{x_1}=t^{g_1-g_2}\frac{y_2}{y_1}=\frac{y_2}{y_1}
\]
uses precisely the hypothesis $g_1=g_2$. Hence each diagonal image is again a graph of the same helicoidal form, with rescaled linear coefficients and pitch, and is minimal. This is the point at which equality of the two weights is essential.

For an independent verification, apply the pair-vector identities of Lemma~\ref{lem:pairvectors} and the dilation identity. The coefficients satisfy
\[
C_{12}=0,
\]
\[
C_{k\ell}=0\quad(3\le k<\ell\le n-1),\qquad
C_{kn}=0\quad(3\le k\le n-1),
\]
and, for $3\le k\le n-1$,
\[
C_{1k}+C_{2k}=-a_k^2(h_{11}+h_{22})=0,
\]
while
\[
C_{1n}+C_{2n}=-(h_{11}+h_{22})=0.
\]
Because $g_1=g_2$, the pairs $\{1,k\}$ and $\{2,k\}$ always lie in the same resonance class. Any additional coincidences among the remaining weights merely merge classes whose displayed coefficient sums are already zero. Hence \eqref{eq:dilationidentity} independently confirms minimality of every diagonal image. Finally,
\[
C_{1n}=-h_{11}=-\frac{2cx_1x_2}{(x_1^2+x_2^2)^2}\ne0
\]
on the chosen quadrant, so \eqref{eq:CijII} proves that the second fundamental form is not identically zero.
\end{proof}

\end{document}
